% !PRINTED 551
\weberpage{551}
% running head: \S.\,151. Ideals. 551
%<<CONT
with the number which in \gesperrt{Dedekind} is called the \gesperrt{norm of the ideal}.
It shall therefore be so named here as well, and accordingly,
when the functional $\varphi$ belongs to the ideal $\mathfrak{a}$, we set

\weq{2}{N_a(\varphi) = N(\mathfrak{a}).}

The norm of an ideal is thus always a natural number.

If $\mathfrak{p}$ is a prime ideal and $p$ is the natural prime
divisible by $\mathfrak{p}$, then

\weq{3}{N(\mathfrak{p}) = p^f,}

and $f$ is called the \gesperrt{degree} of the prime ideal $\mathfrak{p}$.

In the congruences (\S.\,149, 150), according to these stipulations,
the corresponding ideals may also be regarded
as moduli in place of the functionals.

Finally, let a word be added on \gesperrt{Kummer's}
creation of the ideal numbers\footnote{Cf. \gesperrt{Dirichlet-Dedekind}, Zahlentheorie, \S.\,176.}.

If $\alpha, \beta$ are any two numbers from $\mathfrak{o}$ which have the greatest
common divisor $\chi$, and if $\beta = \chi\varphi$, then
the congruence

\weq{4}{\alpha\omega \equiv 0 \pmod{\beta}}

will be satisfied only for those numbers $\omega$ from $\mathfrak{o}$ that are
divisible by $\varphi$. But the functional $\varphi$ will in general
not be associated with any number. Now instead of using the functionals,
one can create for oneself a new system of concepts, in which
one combines all the numbers of $\mathfrak{o}$ into pairs
$(\alpha, \beta)$, and can call these pairs ``\gesperrt{ideal numbers}'' in $\mathfrak{o}$;
in exactly the same way in which, by pairing the integral
rational numbers, one has formed the fractions, or, by pairing the real
numbers, the complex numbers.

One then calls the number $\omega$ divisible by the ideal number $(\alpha, \beta)$
if it satisfies the congruence (1). % [sic] printed reference reads \textquotedbl{}(1)\textquotedbl{}; the congruence introduced on this page is numbered (4)
This is the standpoint
of \gesperrt{Kummer's} theory. But in order that this definition
of divisibility be fruitful, it must still be established
under what assumptions one is to regard two ideal numbers with
different elements, $(\alpha, \beta), (\alpha', \beta')$, as equal,
how the actual numbers in $\mathfrak{o}$ are to be ranged among the ideal
numbers, and finally what is to be understood by the product of two ideal
numbers, so that a number divisible by an ideal number
%CONT>>
